% r2_tui developer's manual
% Build: pdflatex r2_tui_manual.tex (run it twice for the table of contents)
\documentclass[11pt,a4paper]{article}

\usepackage[T1]{fontenc}
\usepackage{lmodern}
\usepackage[margin=2.5cm]{geometry}
\usepackage{xcolor}
\usepackage{listings}
\usepackage{booktabs}
\usepackage{longtable}
\usepackage{array}
\usepackage{tabularx}
\usepackage{enumitem}
\usepackage[hidelinks]{hyperref}

\definecolor{codebg}{gray}{0.96}
\definecolor{comment}{rgb}{0.35,0.45,0.35}
\lstset{
  language=C,
  basicstyle=\ttfamily\small,
  backgroundcolor=\color{codebg},
  commentstyle=\color{comment}\itshape,
  keywordstyle=\bfseries,
  columns=fullflexible,
  keepspaces=true,
  showstringspaces=false,
  frame=single,
  rulecolor=\color{codebg},
  framesep=4pt,
  xleftmargin=6pt,
  xrightmargin=6pt,
  breaklines=true,
  aboveskip=8pt,
  belowskip=8pt
}

% a function, type or macro name (underscores print as they are)
\newcommand{\fn}[1]{\texttt{\detokenize{#1}}}
% a key or key sequence
\newcommand{\key}[1]{\texttt{#1}}
\newcommand{\bs}{\textbackslash}
\newcommand{\prefix}{\key{Ctrl+\bs}}

\setlist{itemsep=2pt,topsep=4pt}
\setlength{\parskip}{4pt}
\setlength{\parindent}{0pt}

\title{\textbf{r2\_tui}\\[4pt]\large Developer's Manual}
\author{Terminal UI for Chestnut and anything else with a terminal}
\date{\texttt{r2\_tui.h} v0.7 \quad \texttt{r2\_tui\_edit.h} v0.4}

\begin{document}

\maketitle

\begin{abstract}
\noindent r2\_tui is a keyboard-driven text user interface in the style of the
Borland IDEs: overlapping and docked windows with shadows, a menu bar that
follows the focused window like the Mac's, message and input boxes, a fuzzy
command palette, widgets and a multi-line text editor. It is a set of
single-file C11 headers that never allocate memory, never block, use no C
library, and talk to the terminal only through bytes you move. This manual is
for developers who build applications on it or port it to a new platform.
\end{abstract}

\tableofcontents
\newpage

%======================================================================
\section{Introduction}

r2\_tui was written as the user interface for Chestnut, a cooperative
operating system that runs on the Raspberry Pi and the ESP32-S3. Chestnut's
graphics library was too slow under its cooperative scheduler, so the UI lives
in a terminal instead: over a serial line to a host terminal, in a Mac or Linux
terminal, or on a device console with a VT100 emulator.

\subsection{Design rules}

Everything in the library follows from a few rules:

\begin{description}[style=nextline]
\item[No allocation.] All memory is given to the library: one arena for the
  screen and window manager, and separate buffers for editor text, clipboard
  and undo history. Sizes are fixed at compile time or by the size of the
  memory you hand over.
\item[Never block.] You feed in the bytes that arrive from the keyboard and pull
  out the bytes to send, in chunks as small as 64 bytes, between which you can
  yield. Nothing waits for input or output.
\item[No C library.] The headers need only \fn{<stdint.h>} and \fn{<stddef.h>}.
  \texttt{make wasm} proves it: the UI and editor link to a WebAssembly module
  with zero imports.
\item[Only changed cells are sent.] A diff renderer compares what the
  application drew with what the terminal shows. Changing one cell costs under
  50 bytes; drawing the same thing again costs nothing.
\item[Keyboard only.] Every feature is reachable without a mouse, without
  \key{Alt}/\key{Option}, and without function keys (which a Mac terminal often
  swallows).
\item[Retained windows, immediate-mode contents.] The window manager keeps
  windows, positions, focus and docking. What is \emph{in} a window is drawn by
  a function you write, every time it runs, without building a widget tree.
\end{description}

\subsection{Files}

\begin{longtable}{@{}p{0.27\textwidth}p{0.68\textwidth}@{}}
\toprule
\textbf{File} & \textbf{What it is} \\
\midrule
\endhead
\fn{r2_tui.h} & The library: cells, width table, drawing, renderer, input
  parser, window manager, menus, message and input boxes, command prefix,
  command palette, widgets. \\
\fn{r2_tui_edit.h} & Optional text editor widget (gap buffer, selection,
  clipboard, undo, search and replace). Must be implemented in the same C file
  as \fn{r2_tui.h}. \\
\fn{r2_tui_posix.h} & Mac and Linux terminal plumbing: raw mode, alternate
  screen, window size, non-blocking reads, a clock. Not used on Chestnut. \\
\fn{r2_termui.h} & VT100 and ANSI escape sequences as macros. \\
\fn{r2_strings.h} & UTF-8 strings (\fn{s8}) used for all text. \\
\fn{examples/tui_app.c} & A starter application: windows, menus, docking,
  dialogs, no editor. \texttt{make app}. \\
\fn{examples/tui_demo.c} & Everything: editor, widgets, test windows.
  \texttt{make demo}. \\
\fn{tools/gen_tui_width.py} & Generates the character width table from
  Python's Unicode data. \\
\fn{tools/wasm_check.c} & The freestanding wasm build check. \texttt{make wasm}. \\
\fn{tests/r2_tui.c}, \fn{tests/r2_tui_edit.c} & The test suites. \\
\bottomrule
\end{longtable}

\subsection{Status}

Everything described here is implemented and tested, with two exceptions
noted where they come up: the editor's syntax highlighting hook
(\fn{R2_TUI_EDIT_NO_HIGHLIGHT} is reserved for it) and the Chestnut window
server (section~\ref{sec:server}). Lisp bindings, so that tool calls from a
language model can drive the UI, are planned.

%======================================================================
\section{Quick start}

\subsection{A complete program}

This is \fn{examples/tui_app.c} cut down to its bones. It opens one window
with a File menu and runs until \key{Ctrl+Q}.

\begin{lstlisting}
#define _DEFAULT_SOURCE        /* POSIX + SIGWINCH for r2_tui_posix.h */
#define _DARWIN_C_SOURCE
#include "r2_tui_posix.h"

#define R2_TUI_IMPLEMENTATION  /* in exactly one C file */
#include "r2_tui.h"

static uint8_t arena[R2_TUI_ARENA_SIZE(400, 200)]; /* largest screen */
static tui t;
static int running = 1;

static void hello(tui *t, tui_window *w, void *user)
{
    tui_menu_begin(t, "~F~ile");
    if (tui_menu_item(t, "E~x~it", "Ctrl+Q"))
        running = 0;
    tui_menu_end(t);
    tui_text(t, 1, 1, S8("Hello from r2_tui"), TUI_DEFAULT);
}

int main(void)
{
    tui_posix p;
    uint8_t in[256], out[1024];
    int cols, rows;
    size_t n;

    tui_init(&t, arena, sizeof(arena), 80, 25);
    tui_posix_open(&p);
    if (tui_posix_size(&cols, &rows) == 0)
        tui_resize(&t, cols, rows);
    tui_window_open(&t, "Hello", (tui_rect){-1, -1, 40, 8}, 0, hello, 0);

    while (running) {
        uint32_t now = tui_posix_now_ms();
        n = tui_posix_read(in, sizeof(in));         /* never blocks */
        tui_feed(&t, in, n, now);                   /* bytes -> keys */
        tui_update(&t, now);                        /* keys -> screen */
        while ((n = tui_render(&t, out, sizeof(out))) > 0)
            tui_posix_write(out, n);                /* changed cells */
        tui_posix_wait(tui_input_pending(&t) ? (int)t.esc_ms : 1000);
    }
    tui_posix_close(&p);
    return 0;
}
\end{lstlisting}

\subsection{Building}

The headers are stb-style: include them anywhere for the declarations, and
define \fn{R2_TUI_IMPLEMENTATION} (and \fn{R2_TUI_EDIT_IMPLEMENTATION} for the
editor) before including in exactly one C file. \fn{r2_tui.h} includes
\fn{r2_strings.h} and \fn{r2_termui.h}, which must sit next to it, and turns on
\fn{R2_STRINGS_IMPLEMENTATION} for you. All functions are \texttt{static}, so
several C files may each have their own implementation without link clashes.

\begin{longtable}{@{}p{0.28\textwidth}p{0.67\textwidth}@{}}
\toprule
\textbf{Make target} & \textbf{Does} \\
\midrule
\endhead
\texttt{make app} & The starter app, \fn{bin/tui_app}. \\
\texttt{make demo} & The full demo, \fn{bin/tui_demo}
  (\texttt{--baud 115200} throttles output to serial speed). \\
\texttt{make test} & All suites with gcc (\texttt{./bin/run\_tests tui} or
  \texttt{edit} runs one). \\
\texttt{make test\_tui}, \texttt{make test\_edit} & One suite. \\
\texttt{make test\_edit\_flags} & The editor suite with each feature
  compiled out. \\
\texttt{make test\_clang} & All suites with clang. \\
\texttt{make wasm} & The freestanding wasm32 check (clang and wasm-ld). \\
\texttt{make check} & cppcheck. \\
\bottomrule
\end{longtable}

The code is C11 and builds without warnings under
\texttt{-Wall -Wextra -Wpedantic -Wshadow} with gcc and clang, and for the
ESP32-S3 with \texttt{xtensa-esp32s3-elf-gcc -mabi=call0}.

%======================================================================
\section{Architecture}

\subsection{Data flow}

\begin{lstlisting}[language={},frame=single]
 keyboard bytes                                      terminal bytes
      |                                                    ^
 tui_feed() ---> input parser ---> event queue             |
                                       |              tui_render()
                                  tui_update()       (diff, chunks)
                                       |                    ^
         route each key: palette, menus, move mode,         |
         modal box, prefix, shortcuts, focused window   back buffer
                                       |                    ^
         run the controllers that need it ("settle")        |
                                       |                    |
         controllers record drawing commands  --->  compositor paints
         into a shared command pool                desktop, docks,
                                                   windows (z-order),
                                                   menus, palette, status
\end{lstlisting}

\begin{enumerate}
\item \fn{tui_feed} turns raw bytes into key events. Escape sequences may be
  split across calls; a lone \key{Esc} waits \fn{esc_ms} (50~ms by default) for
  the rest of a sequence before it counts as the \key{Esc} key.
\item \fn{tui_update} takes the events one at a time. Each goes to whoever
  owns it at that moment (section~\ref{sec:routing}). After \emph{every}
  event the library runs any controller that needs it, so keys that arrive in
  one batch (a serial line, a paste, typing ahead) behave exactly like keys
  typed one by one.
\item A \emph{controller} is your function for a window. It draws the window's
  contents and declares its menus. Its drawing is recorded as commands in a
  shared pool rather than drawn straight away.
\item The compositor replays the recorded commands back to front into the
  back buffer, with frames, shadows, the menu bar, the palette and the status
  line. Moving a window or opening a menu re-composes without running any
  controller.
\item \fn{tui_render} compares the back buffer with the front buffer (what
  the terminal shows) and writes escape sequences for the cells that differ.
\end{enumerate}

\subsection{When controllers run}

A controller runs only when its window needs it:

\begin{itemize}
\item it was just opened, or its inside changed size (resize, dock, zoom);
\item it gained or lost focus (widgets look different when focused);
\item a key was routed to it, or one of its menu items was chosen;
\item once more after that key or menu choice, so that menus and status
  declared \emph{before} the change show the new state;
\item an answer arrived for a message or input box it asked;
\item it was told it may ask again after a busy message box closed;
\item the application called \fn{tui_window_invalidate} (its data changed);
\item the theme or the screen size changed.
\end{itemize}

There is no frame loop: an idle UI runs nothing and sends nothing.

\subsection{Why commands are recorded}

Recording drawing as commands, instead of drawing directly, does two jobs.
Windows can be moved, raised, overlapped and composed under menus without
asking their owners to draw again, which matters when a controller is
expensive (Lisp, for example). And it is the seam for a future window server
on Chestnut (section~\ref{sec:server}): an application task can build its
command list and send it to the task that owns the terminal.

%======================================================================
\section{Integrating with a platform}

\subsection{The loop}

Every platform runs the same five steps; only where the bytes come from and
go to changes.

\begin{lstlisting}
n = read_some_bytes(in, sizeof(in));    /* 0 is fine */
tui_feed(&t, in, n, now_ms);
tui_update(&t, now_ms);                 /* returns 1 if the screen changed */
while ((n = tui_render(&t, out, sizeof(out))) > 0)
    send_bytes(out, n);
sleep_until_input_or(tui_input_pending(&t) ? t.esc_ms : idle_ms);
\end{lstlisting}

\fn{now_ms} is any millisecond counter; only differences are used, so it may
wrap. Call \fn{tui_step} (or \fn{tui_update}) when idle while
\fn{tui_input_pending} is true, so a lone \key{Esc} turns into the \key{Esc}
key on time.

\subsection{POSIX terminals (Mac and Linux)}

\fn{r2_tui_posix.h} is a header of static functions. Define
\fn{_DEFAULT_SOURCE} and \fn{_DARWIN_C_SOURCE} before the first include.

\begin{longtable}{@{}p{0.38\textwidth}p{0.57\textwidth}@{}}
\toprule
\textbf{Function} & \textbf{Does} \\
\midrule
\endhead
\fn{int tui_posix_open(tui_posix *p)} & Raw mode (no echo, no signals, no
  flow control), alternate screen, line wrap off, \fn{SIGWINCH} handler.
  0 on success. \\
\fn{void tui_posix_close(tui_posix *p)} & Restores attributes, cursor, the
  main screen and cooked mode. Always call it. \\
\fn{int tui_posix_size(int *cols, int *rows)} & Size from
  \fn{ioctl(TIOCGWINSZ)}. \\
\fn{int tui_posix_resized(void)} & 1 once after each \fn{SIGWINCH}. \\
\fn{size_t tui_posix_read(void *buf, size_t cap)} & Whatever is waiting;
  never blocks. \\
\fn{int tui_posix_write(const void *buf, size_t n)} & Writes everything. \\
\fn{void tui_posix_wait(int ms)} & Sleeps until input or a timeout
  (\fn{poll}). \\
\fn{uint32_t tui_posix_now_ms(void)} & Monotonic milliseconds. \\
\bottomrule
\end{longtable}

Because raw mode turns signals off, \key{Ctrl+C}, \key{Ctrl+Z} and
\key{Ctrl+\bs} reach the application as keys. The application decides what
quits.

\subsection{Chestnut and serial lines}

On a cooperative OS the UI is a task that must never hold the processor for
long. The renderer is resumable: give it a small buffer, send the chunk,
yield, and call it again; a frame may span many turns. Drawing between chunks
is allowed; the terminal always ends up showing the latest state.

\begin{lstlisting}
for (;;) {                                   /* the UI task */
    uint32_t now = ms();
    size_t n = uart_read(in, sizeof(in));    /* non-blocking */
    tui_feed(&t, in, n, now);
    tui_update(&t, now);
    n = tui_render(&t, out, 128);            /* one chunk */
    if (n)
        uart_write(out, n);
    yield();
}
\end{lstlisting}

At 115200 baud a full 80 by 25 repaint takes over a second, so the diff
renderer matters: ordinary edits cost tens of bytes.

There is no \fn{ioctl} on a serial line. \fn{tui_query_size} writes a request
(save cursor, move to row and column 999, ask for the cursor position, restore)
into your buffer; send it, and the terminal's reply arrives as a
\fn{TUI_EV_RESIZE} event that \fn{tui_update} applies. Over a slow or bursty
line, raise \fn{t.esc_ms} (or \fn{R2_TUI_ESC_MS}) so escape sequences that
arrive in pieces are not mistaken for a lone \key{Esc}.

On the ESP32-S3 (\texttt{-mabi=call0 -Os}) the whole UI with widgets is about
24~KB of code and 4~KB of read-only data; the editor adds about 7~KB, or about
3.7~KB with \fn{R2_TUI_EDIT_MINIMAL}. The 80 by 25 arena is 29~KB and the
\fn{tui} structure about 3~KB.

\subsection{WebAssembly}

\fn{tools/wasm_check.c} builds the UI, the editor, \fn{r2_strings.h} and
\fn{r2_quant.h} for \texttt{wasm32} with clang alone (no emcc, no C library).
wasm-ld fails if anything needs a function the module does not define. The
module exports a small API a JavaScript terminal could drive:
\fn{r2_init}, \fn{r2_in}, \fn{r2_out}, \fn{r2_feed}, \fn{r2_update},
\fn{r2_resize} and \fn{r2_render}. Like every freestanding target it supplies
\fn{memset}, \fn{memcpy}, \fn{memmove} and \fn{memcmp}, which compilers may
call for you (the headers never do).

%======================================================================
\section{Memory and configuration}

\subsection{The arena}

\fn{tui_init(&t, arena, size, cols, rows)} carves everything it needs from
the arena: the back and front cell buffers, combining-mark tables, the window
slots, the shared command pool and two menu bars. \fn{R2_TUI_ARENA_SIZE(cols,
rows)} gives the size needed for a screen of up to that many cells; it is a
constant expression, so the arena can be a static array.

\begin{lstlisting}
static uint8_t arena[R2_TUI_ARENA_SIZE(132, 50)];
tui t;
if (tui_init(&t, arena, sizeof(arena), 80, 25) != 0) /* too small */ ...
tui_resize(&t, 132, 50);  /* fine: fits the arena; clears the screen */
\end{lstlisting}

The \fn{tui} structure itself (about 3~KB) holds the input queue, the parser,
the widget state pool and the palette's results. On a 64-bit machine
\fn{R2_TUI_ARENA_SIZE(80, 25)} is 29\,288 bytes and
\fn{R2_TUI_ARENA_SIZE(132, 50)} is 67\,240.

\subsection{Configuration macros}

Define any of these before including the header (the same value in every file
that includes it).

\begin{longtable}{@{}p{0.33\textwidth}p{0.12\textwidth}p{0.48\textwidth}@{}}
\toprule
\textbf{Macro} & \textbf{Default} & \textbf{Meaning} \\
\midrule
\endhead
\fn{R2_TUI_MAX_WINDOWS} & 8 & User windows open at once (at most 30). Two
  more slots are internal: the desktop and the message box. \\
\fn{R2_TUI_CMD_POOL} & 6144 & Bytes of recorded drawing shared by all
  windows. An editor filling a large window wants 32~KB. \\
\fn{R2_TUI_MAX_COMBINING} & 64 & Combining marks the screen can hold. \\
\fn{R2_TUI_EVENT_QUEUE} & 32 & Keys waiting for \fn{tui_update}. \\
\fn{R2_TUI_ESC_MS} & 50 & How long a lone \key{Esc} waits. \\
\fn{R2_TUI_TITLE_MAX} & 32 & Window title bytes. \\
\fn{R2_TUI_STATUS_MAX} & 96 & Status line bytes per window. \\
\fn{R2_TUI_MENU_MAX} & 8 & Menus per menu bar. \\
\fn{R2_TUI_MENU_ITEMS} & 48 & Items per menu bar. \\
\fn{R2_TUI_MENU_TEXT} & 768 & Label and shortcut bytes per menu bar. \\
\fn{R2_TUI_MESSAGE_MAX} & 256 & Message box text bytes. \\
\fn{R2_TUI_INPUT_MAX} & 128 & Longest text an input box edits. \\
\fn{R2_TUI_PREFIX} & \texttt{"Ctrl+\bs\bs"} & The window command prefix
  (a C string). \\
\fn{R2_TUI_PREFIX_MS} & 2000 & How long the prefix waits for its key. \\
\fn{R2_TUI_PALETTE_KEY} & \texttt{"Ctrl+]"} & The command palette key. \\
\fn{R2_TUI_PALETTE_MAX} & 128 & Palette results. \\
\fn{R2_TUI_MAX_WIDGETS} & 64 & Focusable widgets per window (more draw,
  but Tab skips them). \\
\fn{R2_TUI_STATE_SLOTS} & 32 & Widget state kept between runs (least
  recently used goes). \\
\fn{R2_TUI_EDIT_QUERY} & 64 & Editor search and replace text bytes. \\
\bottomrule
\end{longtable}

When a fixed area runs out the library degrades rather than fails: a window
whose drawing overflows the command pool shows a red \texttt{!} on its frame;
menu declarations past the limits are not recorded; extra combining marks and
events are dropped.

%======================================================================
\section{Screen, drawing and input}

These are the low-level layer. Applications using the window manager rarely
call the \fn{tui_draw_*} functions themselves; windows use \fn{tui_text} and
friends (section~\ref{sec:controllers}).

\subsection{Cells and styles}

A cell (\fn{tui_cell}) is 32 bits: a 21-bit rune and an 11-bit style. A style
(\fn{tui_style}) is a foreground colour (4 bits), a background colour
(4 bits) and attributes (3 bits).

\begin{lstlisting}
tui_style s = TUI_STYLE(TUI_BRIGHT_YELLOW, TUI_BLUE, TUI_BOLD);
\end{lstlisting}

The colours are the 16 ANSI colours, \fn{TUI_BLACK} to \fn{TUI_WHITE} and
\fn{TUI_DARK_GRAY} to \fn{TUI_BRIGHT_WHITE}. Attributes are \fn{TUI_BOLD},
\fn{TUI_UNDERLINE} and \fn{TUI_REVERSE}. \fn{TUI_DEFAULT} is a special style
meaning ``the window's own colours'', resolved when the window is composed.

\subsection{Text}

All text is UTF-8, passed as \fn{s8} (pointer and length, from
\fn{r2_strings.h}): \fn{S8("literal")} or \fn{s8_from_cstr(p)}.

\begin{itemize}
\item Wide characters (CJK, most emoji) take two cells. The right half is the
  rune \fn{TUI_RUNE_CONT}. A wide character cut by the clip or the screen edge
  becomes a space; overwriting half of one blanks the other half.
\item Combining marks attach to the previous cell (one per cell).
\item Joiners, variation selectors, bidirectional controls, the BOM and tag
  characters are dropped, so the terminal cannot merge or reorder glyphs
  behind the layout's back. An emoji ZWJ sequence draws as its parts.
\item Control characters are drawn as \texttt{?}, so text from a file can never
  inject an escape sequence.
\end{itemize}

\fn{tui_rune_width} and \fn{tui_text_width} give columns. The width table is
generated from Unicode 16.0 by \fn{tools/gen_tui_width.py}: about 2.3~KB of
packed ranges, searched by bisection. A test checks the C lookup against the
generator's checksum for all 1.1 million codepoints.

\subsection{Drawing}

\begin{longtable}{@{}p{0.48\textwidth}p{0.47\textwidth}@{}}
\toprule
\textbf{Function} & \textbf{Does} \\
\midrule
\endhead
\fn{tui_draw_rune(t, x, y, r, style)} & One rune; returns columns (0--2). \\
\fn{tui_draw_text(t, x, y, s8, style)} & A line of text; returns columns. \\
\fn{tui_draw_fill(t, rect, rune, style)} & Fill a rectangle. \\
\fn{tui_draw_box(t, rect, style, kind)} & \fn{TUI_BOX_SINGLE} or
  \fn{TUI_BOX_DOUBLE} outline. \\
\fn{tui_draw_shadow(t, rect)} & Darken cells, keeping their characters. \\
\fn{tui_clip(t, rect)}, \fn{tui_clip_reset(t)} & Limit drawing. \\
\fn{tui_set_cursor(t, x, y, visible)} & Where the cursor goes after a frame. \\
\fn{tui_get_cell}, \fn{tui_get_mark} & Read the back buffer (tests). \\
\bottomrule
\end{longtable}

\subsection{The renderer}
\label{sec:renderer}

\fn{size_t tui_render(tui *t, uint8_t *out, size_t cap)} writes the next part
of the frame (at most \fn{cap} bytes, at least \fn{R2_TUI_RENDER_MIN}, 64) and
returns its length, or 0 when the terminal is up to date. Its guarantees:

\begin{itemize}
\item Only differing cells are sent, with the shortest cursor move (none,
  forward, or absolute) and a colour change only when the colour changes.
\item The front buffer is updated as bytes are produced, so an interrupted
  frame stays consistent.
\item A frame is wrapped in synchronized-update brackets
  (\texttt{ESC[?2026h} \ldots\ \texttt{ESC[?2026l}) and hides the cursor while
  drawing. If nothing really changed (something was drawn again identically),
  nothing at all is sent.
\item The renderer never writes past the right edge and never relies on line
  wrap.
\end{itemize}

\fn{tui_render_pending} tells whether a frame is in progress or due.
\fn{tui_invalidate} forgets what the terminal shows and repaints everything
(use it after something else wrote to the terminal).

\subsection{Input}

\fn{tui_feed(t, bytes, n, now_ms)} parses; \fn{tui_poll(t, &ev)} takes the next
event (applications using \fn{tui_update} do not call it). An event
(\fn{tui_event}) has a \fn{type} (\fn{TUI_EV_KEY} or \fn{TUI_EV_RESIZE}), a
\fn{key}, a \fn{ch} (for \fn{TUI_KEY_RUNE}) and \fn{mods}
(\fn{TUI_MOD_SHIFT}, \fn{TUI_MOD_ALT}, \fn{TUI_MOD_CTRL}).

\begin{itemize}
\item Keys: \fn{TUI_KEY_RUNE} (a character), \fn{ENTER}, \fn{TAB},
  \fn{BACKSPACE}, \fn{ESC}, the arrows, \fn{HOME}, \fn{END}, \fn{PGUP},
  \fn{PGDN}, \fn{INSERT}, \fn{DELETE}, \fn{TUI_KEY_F1} to \fn{TUI_KEY_F12}.
\item CSI and SS3 sequences from xterm, VT100 and the Linux console, with
  modifiers (\texttt{ESC[1;5A} is \key{Ctrl+Up}).
\item \key{Ctrl}+letter arrives as \fn{TUI_KEY_RUNE} with the lower-case letter
  and \fn{TUI_MOD_CTRL}. Byte \texttt{0x00} is \key{Ctrl+Space}; \texttt{0x1C}
  to \texttt{0x1F} are \key{Ctrl+\bs}, \key{Ctrl+]}, \key{Ctrl+\textasciicircum}
  and \key{Ctrl+\_} (which is also what \key{Ctrl+/} sends).
  \texttt{0x08} and \texttt{0x7F} are \key{Backspace}; CR and LF are
  \key{Enter}. Terminals cannot report \key{Ctrl+Shift}+letter separately.
\item \key{Esc} followed by a key is \key{Alt}+key (how terminals send
  \key{Alt}, and \key{Option} on a Mac with ``Use Option as Meta key'').
\item A cursor-position reply is a \fn{TUI_EV_RESIZE} only while a
  \fn{tui_query_size} request is outstanding; otherwise the same bytes are
  \key{F3} with modifiers, as xterm intends.
\end{itemize}

\fn{tui_key_name} writes a readable name (\texttt{Ctrl+Alt+F5}, \texttt{Shift+Tab}).

%======================================================================
\section{The window manager}

\subsection{Windows}

\begin{lstlisting}
tui_window *w = tui_window_open(&t, "Title", rect, flags, controller, user);
\end{lstlisting}

The rectangle includes the frame. A width or height of 0 or less asks for the
default (50 by 14, shrunk to fit); \texttt{x} or \texttt{y} below 0 asks for
the next cascade position; dialogs are centred. The result is 0 when all
\fn{R2_TUI_MAX_WINDOWS} slots are in use. Windows get the lowest free number
from 1 to 9, shown on the frame and used by \prefix\ then a digit.

\begin{longtable}{@{}p{0.3\textwidth}p{0.65\textwidth}@{}}
\toprule
\textbf{Flag} & \textbf{Meaning} \\
\midrule
\endhead
\fn{TUI_WIN_DIALOG} & Dialog colours; opened centred; kept inside the
  desktop rather than the work area. \\
\fn{TUI_WIN_MODAL} & Takes all input until closed: no menus, no switching. \\
\fn{TUI_WIN_NO_SHADOW} & No shadow. \\
\fn{TUI_WIN_NO_RESIZE} & Move only; no zoom. \\
\fn{TUI_WIN_ASK_CLOSE} & Closing asks the controller (\fn{tui_closing}
  returns 1 once) instead of closing. \\
\fn{TUI_WIN_NO_CLOSE} & Cannot be closed by the user. \\
\bottomrule
\end{longtable}

Other calls: \fn{tui_window_close}, \fn{tui_window_focus} (also brings back a
minimized window), \fn{tui_window_invalidate} (run its controller),
\fn{tui_window_set_title}, \fn{tui_focused_window},
\fn{tui_window_dock(t, w, side, size)} with \fn{TUI_DOCK_LEFT},
\fn{TUI_DOCK_RIGHT}, \fn{TUI_DOCK_BOTTOM} or \fn{TUI_DOCK_NONE} to float again.
\fn{tui_window} fields such as \fn{frame}, \fn{flags}, \fn{dock},
\fn{zoomed}, \fn{hidden} (minimized) and \fn{runs} may be read; change windows
only through the functions.

\subsection{Layout}

Row 0 is the menu bar and the last row the status line; between them is the
desktop. Docked windows take the left and right edges (full height, 20
columns by default, at most a third of the width) and the bottom between them
(6 rows by default, at most half the height); several on one side share it.
What is left is the \emph{work area} (\fn{tui_work_area}). Floating windows
are kept inside it, shrunk if need be; a zoomed (maximized) window fills it.
A minimized window is not drawn and a minimized dock gives its space back.

\subsection{Focus, z-order and the desktop}

The focused window has a double frame and is on top of the floating windows
(modal windows are always above the rest). \fn{tui_set_desktop(t, controller,
user)} gives the desktop a controller: its menus show when the focused window
declares none, and its status text shows when no window has focus.

\subsection{Routing a key}
\label{sec:routing}

Each key goes to the first of these that wants it:

\begin{enumerate}
\item the command palette, if open;
\item the open menu;
\item move and size mode;
\item the topmost modal window (nothing below gets keys);
\item the key after the command prefix;
\item the prefix itself, then the palette key;
\item the window manager's own keys (Appendix~\ref{app:keys});
\item a shortcut of a menu item on the menu bar;
\item the focused window's controller (or the desktop's).
\end{enumerate}

So menu shortcuts win over a window's own keys. To show a key in a menu
without binding it, write the shortcut in a form the parser does not accept,
such as Emacs notation: \texttt{"C-w"} is displayed but not bound. The demo's
editor uses this so that \key{Ctrl+W} reaches the editor.

\subsection{The command prefix and palette}

\fn{tui_set_prefix(t, "Ctrl+\bs\bs")} sets the window command prefix (0
turns it off; the hints then fall back to function keys).
\fn{tui_set_palette_key(t, "Ctrl+]")} sets the palette key (0 for none;
\prefix\ then \key{/} still opens it). \fn{tui_palette(t)} opens it from code,
including from a menu item.

The palette lists every enabled item of the focused window's and the
desktop's menus (\texttt{File > Save}; a command both declare is listed once),
the window commands, and the other windows. Matching is fuzzy: the query's
letters must appear in order; starts of words and runs of consecutive letters
score higher, gaps and length lower. Results are kept in score order.

\subsection{Themes}

\fn{tui_set_theme(t, &theme)} replaces the colours (0 for the default Borland
theme) and re-runs every window. A \fn{tui_theme} has: \fn{desktop} and
\fn{desktop_rune}; \fn{menu}, \fn{menu_hot}, \fn{menu_sel},
\fn{menu_sel_hot}, \fn{menu_disabled}; \fn{frame}, \fn{frame_active},
\fn{frame_move}; \fn{body}; \fn{dialog_frame}, \fn{dialog_body}; \fn{status},
\fn{status_hot}; \fn{button}, \fn{button_focus}, \fn{button_hot}; \fn{hot},
\fn{focus} (foreground only, on the window's background); \fn{input},
\fn{input_focus}; \fn{list}, \fn{list_sel}, \fn{list_sel_focus};
\fn{scrollbar}.

%======================================================================
\section{Writing controllers}
\label{sec:controllers}

\subsection{The shape of a controller}

\begin{lstlisting}
static void notes(tui *t, tui_window *w, void *user)
{
    struct notes *n = user;
    int cw, ch;

    file_menu(t);                       /* shared menus: just call */
    tui_menu_begin(t, "~N~ote");
    if (tui_menu_item(t, "~R~ename", "F2"))
        n->renaming = 1;
    tui_menu_end(t);

    tui_size(t, &cw, &ch);              /* inside of the window */
    tui_text(t, 1, 0, s8_from_cstr(n->text), TUI_DEFAULT);
    tui_status(t, " ~F2~ Rename  ~Ctrl+\\~ Commands");

    const tui_event *ev = tui_key(t);   /* a key for this window? */
    if (ev && ev->key == TUI_KEY_RUNE && ev->ch == 'x' && !ev->mods) {
        n->count++;
        tui_use_key(t);                 /* handled */
    }
}
\end{lstlisting}

\begin{itemize}
\item Coordinates are relative to the inside of the window (inside the
  frame). \fn{tui_size} gives its width and height.
\item Drawing: \fn{tui_text}, \fn{tui_fill}, \fn{tui_box} and \fn{tui_cursor}
  (where the terminal cursor goes while this window has focus). Use
  \fn{TUI_DEFAULT} for the window's own colours.
\item \fn{tui_key} returns the key this run is for, or 0. Call
  \fn{tui_use_key} if you handled it; an unused \key{Tab} or \key{Shift+Tab}
  moves widget focus.
\item \fn{tui_focused} is 1 when this window has focus. \fn{tui_status} sets the
  status line while it has focus; \texttt{\textasciitilde} marks the
  highlighted parts.
\item Keep application state outside the controller (\fn{user}, or statics):
  the controller is called many times and must be quick and must not block.
\end{itemize}

\subsection{Menus}

Menus are declared on every run, between \fn{tui_menu_begin} and
\fn{tui_menu_end}. \fn{tui_menu_item} returns 1 when that item was chosen
(from the menu, a shortcut or the palette); \fn{tui_menu_item_flags} adds
\fn{TUI_ITEM_DISABLED} and \fn{TUI_ITEM_CHECKED}; \fn{tui_menu_separator}
draws a line. Only the focused window's (and the desktop's) declarations are
recorded; other windows' calls return 0. Because declarations are just calls,
menus shared by every window are a function every controller calls.

A chosen item is identified by its position among the items declared in the
run, so the same controller must declare its menu the same way on the run
that handles the choice (it naturally does). Labels mark their hot letter with
\texttt{\textasciitilde}: \texttt{"\textasciitilde F\textasciitilde ile"}.

Shortcut strings are \emph{modifiers} then a key, joined by \texttt{+}:
modifiers \texttt{Ctrl}, \texttt{Alt}, \texttt{Shift} (any case); keys
\texttt{F1}--\texttt{F12}, \texttt{Enter}, \texttt{Tab}, \texttt{Backspace},
\texttt{Esc}, \texttt{Up}, \texttt{Down}, \texttt{Left}, \texttt{Right},
\texttt{Home}, \texttt{End}, \texttt{PgUp}, \texttt{PgDn}, \texttt{Ins},
\texttt{Insert}, \texttt{Del}, \texttt{Delete}, \texttt{Space}, or a single
character. Anything else is shown but not bound.

\subsection{Message and input boxes}

Both are modal boxes owned by the library, asked for from inside a controller
by calling the function on every run while you want the answer:

\begin{lstlisting}
if (n->quitting) {
    unsigned r = tui_message(t, "Quit", "Leave without saving?",
                             TUI_MB_YES | TUI_MB_NO | TUI_MB_CANCEL);
    if (r == TUI_MB_YES) running = 0;
    if (r) n->quitting = 0;        /* 0 means: still asking */
}
if (n->renaming) {
    unsigned r = tui_input_box(t, "Rename", "New ~n~ame:", n->name,
                               sizeof(n->name));
    if (r) n->renaming = 0;        /* TUI_MB_OK: n->name holds it */
}
\end{lstlisting}

\begin{itemize}
\item The result is 0 while the box is up, then the button
  (\fn{TUI_MB_OK}, \fn{YES}, \fn{NO}, \fn{CANCEL}) once.
\item The box is identified by its title and the window that asked. The answer
  goes only to that window; if it runs without collecting it, the answer is
  dropped.
\item There is one box at a time. Another window asking while it is up gets 0
  and is run again when the box is free.
\item The input box copies the text into your buffer only on OK. It edits at
  most \fn{R2_TUI_INPUT_MAX} bytes.
\item Message box keys: \key{Left}, \key{Right}, \key{Tab} move; \key{Enter} or
  \key{Space} choose; the hot letter chooses; \key{Esc} means Cancel (or the
  last button). Text wraps at 46 columns; \texttt{\bs n} breaks a line.
\end{itemize}

\subsection{Closing}

Without \fn{TUI_WIN_ASK_CLOSE}, closing just closes. With it, the controller
sees \fn{tui_closing} return 1 once and decides, for example by asking with
\fn{tui_message} and calling \fn{tui_window_close} on Yes.

%======================================================================
\section{Widgets}

Widgets are immediate mode: call them on every run; they draw themselves and
report what happened.

\subsection{Layout}

A layout cursor places widgets so you do not pass coordinates.
\fn{tui_at(t, x, y)} moves it; each widget goes below the previous one;
\fn{tui_same_line(t)} puts the next widget beside the last; \fn{tui_space(t,
n)} skips lines.

\begin{lstlisting}
tui_at(t, 2, 1);
tui_label(t, "~N~ame:");  tui_same_line(t);
tui_input(t, "~N~ame##name", name, sizeof(name), 22);
tui_checkbox(t, "~W~rap long lines", &wrap);
tui_radio(t, "~S~paces", &indent, 0);  tui_same_line(t);
tui_radio(t, "~T~abs", &indent, 1);
if (tui_list(t, "schemes", items, 8, &sel, 20, 6) & TUI_CHANGED) ...
if (tui_button(t, "~O~K")) ...
\end{lstlisting}

\subsection{Identity, focus and hot letters}

A widget's identity is a hash of its whole label and its window. Two widgets
with the same visible label are told apart with a suffix that is not shown:
\texttt{"Go\#\#1"}, \texttt{"Go\#\#2"} (repeating a label still works, but
then identity depends on call order).

\key{Tab} and \key{Shift+Tab} move focus in call order, wrapping round, unless
the focused widget used the key (the editor types a tab). A label's hot letter
works on its own unless the focused widget takes letters (a text input, a
list or the editor), in which case \key{Alt}+letter does. Focus survives a
widget disappearing: it moves to one that exists.

State kept between runs (text cursors, scroll positions) lives in
\fn{R2_TUI_STATE_SLOTS} slots in the \fn{tui} structure; the least recently
used slot is recycled and starts clean.

\subsection{Reference}

\begin{longtable}{@{}p{0.45\textwidth}p{0.5\textwidth}@{}}
\toprule
\textbf{Widget} & \textbf{Keys and result} \\
\midrule
\endhead
\fn{tui_label(t, text)} & Text with a hot letter; not focusable. \\
\fn{int tui_button(t, label)} & \key{Enter}, \key{Space}, hot letter.
  Returns 1 when pressed. \\
\fn{int tui_checkbox(t, label, &value)} & \key{Space}, hot letter. Returns 1
  when toggled. \\
\fn{int tui_radio(t, label, &value, option)} & \key{Space}, hot letter. Sets
  \fn{*value} to \fn{option}; returns 1 if that changed it. \\
\fn{unsigned tui_input(t, label, buf, cap, w)} & One line of UTF-8 text in
  \fn{buf} (\fn{cap} bytes with the NUL), \fn{w} columns, scrolling sideways.
  Arrows, \key{Home}, \key{End}, \key{Backspace}, \key{Delete},
  \key{Ctrl+A/E/B/F/D/K}. Returns \fn{TUI_CHANGED} and/or \fn{TUI_SUBMIT}
  (\key{Enter}). \\
\fn{unsigned tui_list(t, label, items, count, &sel, w, h)} & Arrows,
  \key{PgUp}, \key{PgDn}, \key{Home}, \key{End}, first letter. Returns
  \fn{TUI_CHANGED} and/or \fn{TUI_SUBMIT}. Scrollbar when needed. \\
\fn{tui_scroll_begin(t, label, w, h, content_h)} \ldots\
  \fn{tui_scroll_end(t)} & A \fn{w} by \fn{h} view of \fn{content_h} lines drawn
  between the calls at \texttt{y = 0 .. content\_h}; clipped and scrolled.
  Focused: arrows, \key{PgUp}, \key{PgDn}, \key{Space}, \key{Home}, \key{End}.
  Views do not nest; widgets inside are not scrolled into view. \\
\bottomrule
\end{longtable}

%======================================================================
\section{The text editor}

\fn{r2_tui_edit.h} adds a multi-line editor as a widget. Implement it in the
same C file as \fn{r2_tui.h} (it uses the widget internals; the header stops
with an error otherwise).

\subsection{Setting up}

\begin{lstlisting}
#define R2_TUI_IMPLEMENTATION
#include "r2_tui.h"
#define R2_TUI_EDIT_IMPLEMENTATION
#include "r2_tui_edit.h"

static char text[16 * 1024], clip_mem[4096], history[8192];
static tui_edit ed;
static tui_clipboard clip;

tui_edit_init(&ed, text, sizeof(text));
tui_clipboard_init(&clip, clip_mem, sizeof(clip_mem));
ed.clip = &clip;                               /* optional */
tui_edit_undo_init(&ed, history, sizeof(history)); /* optional */
ed.flags = TUI_EDIT_LINE_NUMBERS | TUI_EDIT_AUTO_INDENT;
tui_edit_set(&ed, S8("hello\nworld"));

static void editor(tui *t, tui_window *w, void *user)
{
    int cw, ch;
    char line[100];
    tui_size(t, &cw, &ch);
    tui_at(t, 0, 0);
    if (tui_editor(t, "##ed", &ed, cw, ch) & TUI_CHANGED) ...
    if (tui_edit_prompt(&ed, line, sizeof(line)))  /* searching? */
        tui_status(t, line);
}
\end{lstlisting}

\subsection{The editor structure}

\fn{tui_edit} is public so applications can read it: \fn{cursor} (a byte
offset), \fn{line} (from 0), \fn{lines}, \fn{top_line}, \fn{left},
\fn{modified} (set by any change; clear it yourself after saving),
\fn{mark_on} and \fn{mark}, \fn{upos} and \fn{ulen} (something to undo when
\fn{upos} is above 0, to redo when \fn{upos} is below \fn{ulen}). Set
\fn{flags}, \fn{tab} (tab width, 4) and \fn{clip}; change the rest only
through the functions.

\begin{longtable}{@{}p{0.3\textwidth}p{0.65\textwidth}@{}}
\toprule
\textbf{Flag} & \textbf{Meaning} \\
\midrule
\endhead
\fn{TUI_EDIT_LINE_NUMBERS} & A gutter with line numbers. \\
\fn{TUI_EDIT_READ_ONLY} & Moving, selecting and copying only. \\
\fn{TUI_EDIT_AUTO_INDENT} & \key{Enter} copies the line's leading blanks. \\
\fn{TUI_EDIT_SPACES} & \key{Tab} inserts spaces to the next tab stop. \\
\bottomrule
\end{longtable}

\subsection{Functions}

\begin{longtable}{@{}p{0.47\textwidth}p{0.48\textwidth}@{}}
\toprule
\textbf{Function} & \textbf{Does} \\
\midrule
\endhead
\fn{tui_edit_init(e, buf, cap)} & Use \fn{buf} for the text; starts empty. \\
\fn{int tui_edit_set(e, s8)} & Replace the text, cursor at the start, new
  history. 0 if it does not fit. \\
\fn{uint32_t tui_edit_len(e)} & Bytes of text. \\
\fn{uint32_t tui_edit_get(e, out, cap)} & Copy out (NUL terminated). \\
\fn{int tui_edit_insert(e, s8)} & Insert at the cursor. 0 if it does not
  fit (nothing changes). \\
\fn{tui_edit_delete(e, from, to)} & Delete bytes \texttt{[from, to)}. \\
\fn{tui_edit_goto(e, offset)} & Move the cursor (snapped to a character). \\
\fn{uint32_t tui_edit_column(e)} & The cursor's display column. \\
\fn{unsigned tui_editor(t, label, e, w, h)} & The widget. \fn{TUI_CHANGED}
  when the text changed. \\
\fn{tui_clipboard_init(c, buf, cap)} & A clipboard (shareable). \\
\fn{int tui_edit_selection(e, &from, &to)} & The selected range, if any. \\
\fn{tui_edit_copy}, \fn{tui_edit_cut}, \fn{int tui_edit_paste} & For
  menus. \\
\fn{tui_edit_undo_init(e, buf, cap)} & Undo history memory. \\
\fn{int tui_edit_undo}, \fn{int tui_edit_redo} & One step; 1 if done. \\
\fn{tui_edit_undo_group(e)} & Start a new undo step. \\
\fn{tui_edit_search(e, backward)} & Start an incremental search. \\
\fn{int tui_edit_find(e, s8, backward)} & Find from the cursor and select
  (without selection compiled in: the cursor goes past the match). \\
\fn{tui_edit_replace(e)} & Start a query replace. \\
\fn{int tui_edit_prompt(e, buf, cap)} & The line to show while searching or
  replacing; 0 if none. \\
\bottomrule
\end{longtable}

\subsection{How it works}

The text is a gap buffer in your memory. The editor keeps the line count, the
cursor's line and the first line on screen up to date as it edits, so a key
press never rescans the whole text. Invalid UTF-8 steps one byte at a time in
both directions; tabs and wide characters are measured in display columns.

\textbf{Selection.} The region is between \fn{mark} and \fn{cursor}. A region
made with \key{Ctrl+Space} follows Emacs: moving extends it, typing does not
replace it. One made with \key{Shift}+movement follows CUA: a plain move
clears it and typing replaces it.

\textbf{Clipboard.} Kills (\key{Ctrl+W}, \key{Ctrl+K}, \key{Esc d},
\key{Esc Backspace}) and copies go to the clipboard. Kills in a row join up,
appending or prepending by direction. When the clipboard is full the text is
cut at a character boundary.

\textbf{Undo.} History is a log of records, oldest first:
kind (1 byte), position (4), length (4), step number (4), the text, and the
record size (4) so undo can walk backwards. An insert records the text it put
in, a delete the text it took out; typing into the same record grows it.
Every key is a step, except that up to 20 typed characters join one step;
\fn{tui_edit_cut} and \fn{tui_edit_paste} always start their own. When the
memory fills, whole steps are dropped from the old end (never part of a
step); a single change larger than the memory clears the history. A new
change after undo ends redo.

\textbf{Search.} Matching is byte-wise with ASCII case folding when the query
has no capital letters (Emacs' smart case). Query replace asks with two input
boxes; the first starts empty and shows the last search as its default.

\subsection{Feature switches}

Define before including; each removes its code and keys.

\begin{longtable}{@{}p{0.38\textwidth}p{0.57\textwidth}@{}}
\toprule
\textbf{Macro} & \textbf{Removes} \\
\midrule
\endhead
\fn{R2_TUI_EDIT_NO_SELECTION} & The mark and Shift selection. \\
\fn{R2_TUI_EDIT_NO_CLIPBOARD} & Kill, yank, cut, copy, paste. \\
\fn{R2_TUI_EDIT_NO_UNDO} & Undo and redo. \\
\fn{R2_TUI_EDIT_NO_SEARCH} & Incremental search and query replace. \\
\fn{R2_TUI_EDIT_NO_HIGHLIGHT} & Reserved for the highlighting hook (not
  written yet). \\
\fn{R2_TUI_EDIT_MINIMAL} & All of the above. \\
\bottomrule
\end{longtable}

On the ESP32-S3 the editor is about 7.2~KB of code with everything and
3.7~KB minimal.

%======================================================================
\section{Testing}

The tests are plain C, in the style of the rest of r2 (\fn{r2_unit.h}). Each
suite also compiles on its own.

\begin{itemize}
\item \textbf{A model terminal.} The renderer's output is applied byte by byte
  to a small VT100 model (cursor moves, colours, wide-character rules as
  xterm has them) which must end up equal to the back buffer. Thousands of
  random drawing rounds run with random chunk sizes, drawing in the middle of
  frames.
\item \textbf{Keys as bytes.} Window manager, widget and editor tests type real
  key bytes and read the composed screen. A random-key fuzz runs on each
  layer with invariants (focus on an open window, text valid UTF-8, cursor on
  screen).
\item \textbf{Models.} The editor's incremental bookkeeping is checked against
  a plain byte array over 20\,000 random edits; undo and redo against
  snapshots, including with a history too small to keep everything.
\item \textbf{Exhaustive checks.} The width table for every codepoint; the
  key parser for every key re-fed split at every byte.
\item \textbf{Mutation testing.} During development each behaviour was broken
  on purpose to see a test fail; gaps found that way became tests.
\item \textbf{Builds.} gcc and clang, AddressSanitizer and UBSan, every editor
  feature switch, the wasm module, and the ESP32-S3 compiler.
\end{itemize}

Interactive checks in a real terminal are by hand: \texttt{make demo} and
\texttt{make app}.

%======================================================================
\section{Extending}

\subsection{A new widget}

A widget follows the pattern of those in \fn{r2_tui.h}:

\begin{lstlisting}
unsigned my_widget(tui *t, const char *label, ...)
{
    uint32_t id = r2__tui_wid(t, label);           /* identity */
    int focused = r2__tui_wreg(t, id, KIND);       /* focusable */
    tui_wstate *st = r2__tui_state(t, id);         /* kept state */
    const tui_event *ev = tui_key(t);
    int x, y;
    if (r2__tui_hot_hit(t, label))                 /* hot letter */
        r2__tui_wfocus(t, id, KIND);
    if (focused && ev) { /* handle it */ tui_use_key(t); }
    r2__tui_place(t, w, h, &x, &y);                /* layout */
    /* draw with tui_text / tui_fill at (x, y) */
}
\end{lstlisting}

The \fn{r2__} functions are internal: they are visible in the C file that has
the implementation and may change.

\subsection{Other changes}

\begin{itemize}
\item \textbf{Keys.} New escape sequences go in \fn{r2__tui_csi} and
  \fn{r2__tui_ss3}; add them to the key table test.
\item \textbf{The width table.} After a Python upgrade with newer Unicode data,
  run \texttt{python3 tools/gen\_tui\_width.py r2\_tui.h}; it rewrites the
  generated block and the checksum the test compares.
\item \textbf{A theme.} Copy the Borland theme, change colours, pass it to
  \fn{tui_set_theme}.
\end{itemize}

\subsection{The Chestnut window server}
\label{sec:server}

Planned, not built. One Chestnut task would own the terminal, focus and the
menu bar; each application task would own its windows and send its command
lists and menu declarations to it. The recorded command format and the
per-window ownership in this library are what make that split possible.

%======================================================================
\section{Limitations}

\begin{itemize}
\item Keyboard only; no mouse.
\item 16 colours (no 256-colour or true-colour themes).
\item One combining mark per cell; emoji ZWJ sequences draw as their parts.
\item One message or input box at a time.
\item Scroll views do not nest and do not follow widget focus.
\item The editor has one cursor, a fixed-size buffer, and no highlighting yet.
\item Terminal.app cannot report \key{Ctrl+Shift}+letter, \key{Ctrl+Shift+Space}
  or \key{Ctrl+Esc} as separate keys; \key{Alt} needs ``Use Option as Meta
  key''; macOS may take \key{Ctrl+Space} for switching input sources.
\item Over \texttt{idf.py monitor}, \key{Ctrl+]} quits the monitor; set another
  palette key there.
\end{itemize}

%======================================================================
\appendix

\section{Key reference}
\label{app:keys}

\subsection{Window commands (\prefix\ then a key)}

\begin{longtable}{@{}p{0.3\textwidth}p{0.65\textwidth}@{}}
\toprule
\textbf{After the prefix} & \textbf{Does} \\
\midrule
\endhead
a letter & Open the menu with that hot letter (\key{Shift} is fine). \\
\key{1}--\key{9} & Focus that window (and bring it back if minimized). \\
\key{Space} & The system menu. \\
\key{/} & The command palette. \\
\key{Up} & Maximize. \\
\key{Down} & Restore a maximized window, otherwise minimize. \\
\key{Left}, \key{Right}, \key{Tab}, \key{Shift+Tab} & Previous or next
  window. \\
\key{Enter} & Move (arrows) and size (\key{Shift}+arrows); \key{Enter} done,
  \key{Esc} back. A docked window's arrows change its size. \\
\key{Backspace}, \key{Delete} & Close. \\
\prefix & Send the prefix to the window. \\
\key{Esc}, any other key, 2 seconds & Cancel. \\
\bottomrule
\end{longtable}

\subsection{Other window manager keys}

\begin{longtable}{@{}p{0.3\textwidth}p{0.65\textwidth}@{}}
\toprule
\textbf{Key} & \textbf{Does} \\
\midrule
\endhead
\key{Ctrl+]} & The command palette (\key{Up}/\key{Down} or
  \key{Ctrl+P}/\key{Ctrl+N} choose, \key{Enter} runs, \key{Backspace} and
  \key{Ctrl+U} edit, \key{Esc} or \key{Ctrl+G} close). \\
\key{F10}, \key{Alt}+letter & Menus. \\
\key{Alt+Space} & The system menu. \\
\key{F6}, \key{Shift+F6}, \key{Alt+1}--\key{9} & Switch windows. \\
\key{F5} & Zoom. \\
\key{Ctrl+F5} & Move and size. \\
\key{Alt+F3} & Close. \\
In a menu & Arrows move, \key{Enter} or the hot letter chooses, \key{Esc}
  closes. \\
\bottomrule
\end{longtable}

\subsection{Editor}

\begin{longtable}{@{}p{0.36\textwidth}p{0.59\textwidth}@{}}
\toprule
\textbf{Key} & \textbf{Does} \\
\midrule
\endhead
Arrows, \key{Home}, \key{End}, \key{PgUp}, \key{PgDn} & Move. \\
\key{Ctrl+Home}, \key{Ctrl+End} & Top, bottom. \\
\key{Ctrl+Left}, \key{Ctrl+Right} & Words. \\
\key{Ctrl+A}, \key{Ctrl+E} & Start, end of the line. \\
\key{Ctrl+B}, \key{Ctrl+F}, \key{Ctrl+P}, \key{Ctrl+N} & Back, forward, up,
  down. \\
\key{Ctrl+D}, \key{Delete}, \key{Backspace} & Delete. \\
\key{Ctrl+K} & Kill to the end of the line (or join the next). \\
\key{Enter}, \key{Tab} & New line, tab. \key{Shift+Tab} leaves the editor. \\
\key{Esc} then a key & Meta: \key{Esc f}/\key{b} words, \key{Esc <}/\key{>}
  top and bottom, \key{Esc v} page up, \key{Esc d} and \key{Esc Backspace}
  kill a word, \key{Esc w} copy, \key{Esc \%} replace. \key{Esc Esc}
  clears the selection. \\
\key{Ctrl+Space} & Set the mark. \\
\key{Shift}+movement & Select (CUA). \\
\key{Ctrl+G} & Clear the selection; leave a search. \\
\key{Ctrl+W}, \key{Ctrl+Y} & Kill (cut), yank (paste). \\
\key{Shift+Del}, \key{Ctrl+Ins}, \key{Shift+Ins} & Cut, copy, paste
  (Borland). \\
\key{Ctrl+/}, \key{Ctrl+Z} & Undo. \\
\key{Esc} then \key{Ctrl+/} & Redo. \\
\key{Ctrl+S}, \key{Ctrl+R} & Incremental search forward, backward. \\
In a search & Type to search; \key{Ctrl+S}/\key{Ctrl+R} next (again after a
  failure wraps; on an empty search, the last one); \key{Backspace} shorter;
  \key{Enter} or \key{Esc} stay; \key{Ctrl+G} back; other keys stop. \\
In query replace & \key{y} or \key{Space} replace, \key{n} skip, \key{!}
  the rest, \key{.} this one and stop, \key{q}, \key{Enter} or \key{Esc}
  stop. \\
\bottomrule
\end{longtable}

\section{API index}

\subsection{r2\_tui.h}

\begin{description}[style=nextline,font=\normalfont\bfseries]
\item[Setup] \fn{tui_init}, \fn{tui_resize}, \fn{tui_invalidate},
  \fn{R2_TUI_ARENA_SIZE}
\item[Text] \fn{tui_rune_width}, \fn{tui_text_width}
\item[Drawing] \fn{tui_clip}, \fn{tui_clip_reset}, \fn{tui_draw_rune},
  \fn{tui_draw_text}, \fn{tui_draw_fill}, \fn{tui_draw_box},
  \fn{tui_draw_shadow}, \fn{tui_set_cursor}, \fn{tui_get_cell},
  \fn{tui_get_mark}
\item[Output] \fn{tui_render}, \fn{tui_render_pending}
\item[Input] \fn{tui_feed}, \fn{tui_step}, \fn{tui_poll},
  \fn{tui_input_pending}, \fn{tui_query_size}, \fn{tui_key_name}
\item[Windows] \fn{tui_update}, \fn{tui_window_open}, \fn{tui_window_close},
  \fn{tui_window_dock}, \fn{tui_window_focus}, \fn{tui_window_invalidate},
  \fn{tui_window_set_title}, \fn{tui_focused_window}, \fn{tui_set_desktop},
  \fn{tui_set_theme}, \fn{tui_set_prefix}, \fn{tui_set_palette_key},
  \fn{tui_palette}, \fn{tui_work_area}
\item[In a controller] \fn{tui_text}, \fn{tui_fill}, \fn{tui_box},
  \fn{tui_cursor}, \fn{tui_size}, \fn{tui_key}, \fn{tui_use_key},
  \fn{tui_focused}, \fn{tui_status}, \fn{tui_closing}
\item[Menus] \fn{tui_menu_begin}, \fn{tui_menu_end}, \fn{tui_menu_item},
  \fn{tui_menu_item_flags}, \fn{tui_menu_separator}
\item[Boxes] \fn{tui_message}, \fn{tui_input_box}
\item[Widgets] \fn{tui_at}, \fn{tui_same_line}, \fn{tui_space},
  \fn{tui_label}, \fn{tui_button}, \fn{tui_checkbox}, \fn{tui_radio},
  \fn{tui_input}, \fn{tui_list}, \fn{tui_scroll_begin}, \fn{tui_scroll_end}
\end{description}

\subsection{r2\_tui\_edit.h}

\begin{description}[style=nextline,font=\normalfont\bfseries]
\item[Text] \fn{tui_edit_init}, \fn{tui_edit_set}, \fn{tui_edit_len},
  \fn{tui_edit_get}, \fn{tui_edit_insert}, \fn{tui_edit_delete},
  \fn{tui_edit_goto}, \fn{tui_edit_column}
\item[Widget] \fn{tui_editor}
\item[Selection and clipboard] \fn{tui_clipboard_init},
  \fn{tui_edit_selection}, \fn{tui_edit_copy}, \fn{tui_edit_cut},
  \fn{tui_edit_paste}
\item[Undo] \fn{tui_edit_undo_init}, \fn{tui_edit_undo}, \fn{tui_edit_redo},
  \fn{tui_edit_undo_group}
\item[Search] \fn{tui_edit_search}, \fn{tui_edit_find},
  \fn{tui_edit_replace}, \fn{tui_edit_prompt}
\end{description}

\subsection{r2\_tui\_posix.h}

\fn{tui_posix_open}, \fn{tui_posix_close}, \fn{tui_posix_size},
\fn{tui_posix_resized}, \fn{tui_posix_read}, \fn{tui_posix_write},
\fn{tui_posix_write_str}, \fn{tui_posix_wait}, \fn{tui_posix_now_ms}

\section{What goes to the terminal}

The library only ever sends these, so a terminal or emulator (such as the
planned Chestnut console) needs nothing more:

\begin{longtable}{@{}p{0.35\textwidth}p{0.6\textwidth}@{}}
\toprule
\textbf{Bytes} & \textbf{Meaning} \\
\midrule
\endhead
\texttt{ESC[}\textit{row}\texttt{;}\textit{col}\texttt{H} & Move the cursor. \\
\texttt{ESC[}\textit{n}\texttt{C} & Move the cursor right. \\
\texttt{ESC[0;}\ldots\texttt{m} & Colours and attributes: always starting
  with a reset; \texttt{1} bold, \texttt{4} underline, \texttt{7} reverse,
  \texttt{30}--\texttt{37} and \texttt{90}--\texttt{97} foreground,
  \texttt{40}--\texttt{47} and \texttt{100}--\texttt{107} background. \\
\texttt{ESC[?2026h}, \texttt{ESC[?2026l} & Synchronized update on and off
  (ignored by terminals that lack it). \\
\texttt{ESC[?25l}, \texttt{ESC[?25h} & Hide and show the cursor. \\
UTF-8 text & Printable characters only (controls are drawn as \texttt{?}). \\
\texttt{ESC7 ESC[999;999H ESC[6n ESC8} & Size query
  (\fn{tui_query_size}). \\
\texttt{ESC[?1049h}, \texttt{ESC[?7l} & Alternate screen, wrap off (the
  POSIX adapter only). \\
\bottomrule
\end{longtable}

\end{document}
